\bookchapter{Chapter 5: Make Your Phone Understand You Better}

\chapterintro{A few minutes of setup brings a huge boost to your experience. These settings are like “tuning” your phone to better match your usage habits.}

\booksection{I. Sound and Vibration}

Set a ringtone and notification sound you like. Use silent mode (for meetings) and Do Not Disturb mode (for sleeping) flexibly—the world instantly quiets down.

\begin{itemize}
  \item \textbf{Ringtone}: Pick one you can hear and isn't embarrassing.
  \item \textbf{Vibration mode}: Vibration alerts when on silent.
  \item \textbf{Do Not Disturb mode}: Can set exceptions, like allowing calls from contacts.
  \item \textbf{Volume separation}: Media volume, ringtone volume, and alarm volume can be adjusted separately.
\end{itemize}

\booksection{II. Display and Brightness}

\begin{itemize}
  \item \textbf{Auto brightness}: Lets the phone adjust based on ambient light—worry-free.
  \item \textbf{Eye comfort/Night Shift mode}: Filters blue light, making phone use more comfortable at night.
  \item \textbf{Dark mode}: Saves power (especially effective on OLED screens) and isn't harsh on the eyes at night. Many apps also support dark mode.
  \item \textbf{Screen refresh rate}: Adjustable on some Android phones—higher (like 120Hz) means smoother scrolling.
  \item \textbf{Auto-lock}: Set to 15 or 30 seconds—both secure and power-saving.
  \item \textbf{Font size}: Think the text is too small? Make it bigger.
\end{itemize}

\booksection{III. Wallpaper, Themes, and Home Screen}

Change to a nice wallpaper (static or dynamic)—even your mood changes. Android brand phones often have theme stores that can change the entire set of icons and style.

\textbf{Home screen widgets}: Put weather, schedule, and music controls on the home screen for one-step access. Both iOS and Android support this—long-press an empty area on the home screen to add.

\booksection{IV. Efficiency Settings: Make Your Phone Your “External Brain”}

These are “once set, never go back” features—each one can significantly boost your daily efficiency.

\booksubsection{Focus / Do Not Disturb Mode: Reclaim Attention Stolen by Notifications}

\begin{itemize}
  \item \textbf{Core idea}: Not "look at your phone less," but "focus when you should focus, connect when you should connect."
  \item \textbf{Usage suggestions}:\newline · Work mode: Only allow "starred" contacts and essential apps through;\newline · Sleep mode: Auto-enable, only allow alarms;\newline · Fitness / Reading / Meetings: customize based on the current activity.
  \item \textbf{Advanced usage}: Use automation + location / time triggers to switch modes automatically. For example: iOS Shortcuts can set "arrive at company auto-enter work mode."
\end{itemize}

\booksubsection{Notification Management: Stop Red Badges from “Relentlessly Pestering”}

\begin{itemize}
  \item \textbf{Turn off marketing notifications}: 90\% of push notifications you don't actually need.
  \item \textbf{Sort by importance}: Pin interpersonal communications (phone, SMS, instant messaging) to the top; news, shopping, and social platform notifications can be set to "silent delivery."
  \item \textbf{Group notifications}: Fold similar notifications together to avoid being overwhelmed by dozens of messages with one unlock.
\end{itemize}

\booksubsection{Quick Gestures: Stay One Step Ahead}

\begin{itemize}
  \item \textbf{Screenshot}: Power + Volume Down (iOS full-screen) / Power + Volume Down (Android universal).
  \item \textbf{Long screenshot / Scroll screenshot}: After a regular screenshot, tap the "Long Screenshot" option.
  \item \textbf{One-tap scan}: Add a "Scan" icon to Control Center / Quick Settings.
  \item \textbf{Back Tap (iPhone)}: In [Settings] > [Accessibility] > [Touch] > [Back Tap], bind actions like "Screenshot" or "Lock Screen."
  \item \textbf{Voice wake}: Siri / Xiaoai / Xiaoyi / Xiaobu—just call out to check weather or set an alarm.
\end{itemize}

\booksubsection{Input Method: The Most Underestimated Efficiency Tool}

\begin{itemize}
  \item \textbf{Choose the right input method}:\newline · Sogou / Baidu: Large lexicon, fast trending word updates;\newline · iFlytek: Highest voice input recognition rate, a note-taking lifesaver in meetings;\newline · Gboard (Google): Clean and simple, strong cross-platform sync, supports gesture typing;\newline · iOS built-in: Best privacy, deeply integrated with the system.
  \item \textbf{Essential toggles}: Smart correction, word prediction, cloud lexicon sync, swipe up for numbers and symbols.
  \item \textbf{Privacy tip}: Use "full-sentence prediction" or "personalized recommendations" with caution—they may upload your input content.
\end{itemize}

\booksubsection{Home Screen and Multitasking: Smooth Out Your Workflow}

\begin{itemize}
  \item \textbf{App split screen / Picture-in-Picture}: Reply to messages while watching video, take notes while browsing.
  \item \textbf{Small window / Floating window}: WeChat / TikTok support "small window mode"—reply and dismiss.
  \item \textbf{Pin frequently used apps}: Only put the 4 most frequent apps in the Dock (bottom of screen) to reduce search time.
\end{itemize}

\begin{tcbinfo}{"First Day of Work" Checklist}
After getting a new phone or a major version update, spend 10 minutes going through these settings:\newline (1) Lock screen password + biometrics (2) Notification grouping + turn off marketing pushes (3) Font size / refresh rate / dark mode\newline (4) Default input method / browser / email (5) Emergency contacts / medical ID (6) Backup settings

\end{tcbinfo}

\booksection{V. Accessibility: The Power of Tech for Good}

Don't be intimidated by the name “accessibility”—these features are for everyone. They were originally designed for people with disabilities, but they're also useful for ordinary people—temporarily injured hands, holding a baby, broken earbuds on the subway, or needing reading glasses for small text... they can all save the day. Ultimately, accessible design embodies the philosophy of\textbf{tech for good}: letting everyone equally enjoy digital life.

\booksubsection{“Eyes” for Visually Impaired Users—Complete Screen Reader Guide}

Screen readers are the core tool for visually impaired users to use phones. It reads\textbf{all visible content on the screen aloud}, and users operate the phone through gestures and voice. Below we introduce the usage of iOS “VoiceOver” and Android “TalkBack” respectively.

\begin{tcbinfo}{iOS VoiceOver}
\textbf{How to enable}: [Settings] > [Accessibility] > [VoiceOver], turn on the switch. You can also triple-click the side button (or Home button) to quickly toggle.

\textbf{Basic gestures}:

\begin{enumerate}
  \item One-finger swipe: Browse screen content element by element;
  \item One-finger double-tap: Activate element;
  \item :
  \item Three-finger swipe: Switch pages;
  \item One-finger triple-tap: Zoom or select (different from regular triple-tap).
\end{enumerate}

\textbf{Personalization}: [VoiceOver] > [Speech], you can adjust\textbf{speaking rate},\textbf{pitch}, and switch between different voice packs. iOS has built-in “Enhanced Speech” and “Speech Controller”—the former reduces breathing sounds and pauses for a more natural listening experience.

\end{tcbinfo}

\begin{tcbinfo}{Android TalkBack}
\textbf{How to enable}: [Settings] > [Accessibility] > [TalkBack], turn on the switch. Or long-press both volume keys for 3 seconds to quickly enable (supported by some brands).

\textbf{Basic gestures}:

\begin{enumerate}
  \item One-finger left/right swipe: Switch focus element;
  \item One-finger double-tap: Activate element;
  \item Two-finger up/down/left/right swipe: Scroll;
  \item Three-finger single-tap: Enter reading mode (continuous reading);
  \item Four-finger single-tap: Pause/resume reading.
\end{enumerate}

\textbf{Personalization}: [TalkBack] > [Text-to-Speech Settings], you can adjust\textbf{speech rate},\textbf{pitch}, and switch\textbf{the speech engine}(we recommend installing the Google Text-to-Speech engine for more natural reading).

\end{tcbinfo}

\booksubsection{Make Small Text Big, Make Blurry Clear—Visual Adjustments}

\booksubsection{Zoom Gestures}

\begin{itemize}
  \item \textbf{iOS}: [Settings] > [Accessibility] > [Zoom], after enabling,\textbf{three-finger tap the screen}to zoom in. Full-screen zoom is the default, with optional window zoom (a partial magnification pane).\textbf{Pinch with two fingers to adjust magnification (1.5x to 15x)}, three-finger double-tap to lock the current zoom ratio.
  \item \textbf{Android}: [Settings] > [Accessibility] > [Magnification], three-finger tap to zoom, three-finger tap again to exit; or directly\textbf{three-finger triple-tap the screen}to enable temporary magnification (up to 8x).
\end{itemize}

\booksubsection{High Contrast Text / Bold Text}

\begin{itemize}
  \item \textbf{iOS}: [Settings] > [Accessibility] > [Display \& Text Size], turn on "Bold Text," "Increase Contrast," "Reduce Transparency," etc. "Smart Invert" can keep text black on light backgrounds, avoiding white-on-white.
  \item \textbf{Android}: [Settings] > [Display] > [Font Size], with four levels: "Small," "Default," "Large," and "Extra Large." Huawei/Xiaomi additionally provide "Bold" and "High Contrast Text" switches.
\end{itemize}

\booksubsection{Color Filters (Color Blind-Friendly)}

\begin{itemize}
  \item \textbf{iOS}: [Settings] > [Accessibility] > [Display \& Text Size] > [Color Filters], provides\textbf{four color blind modes: red/green/blue/yellow}, plus "Grayscale" and "Invert Colors" options. Take the built-in color vision test once before choosing a filter.
  \item \textbf{Android}: [Settings] > [Accessibility] > [Color Correction], includes\textbf{three color blind modes}plus an "Invert Colors" switch.
\end{itemize}

\booksubsection{AI Image Description: Let AI Be Your “Eyes”}

When a screen reader encounters an image, the traditional approach is to skip it or just read “image.” But now AI can\textbf{describe image content in a sentence}, which is very helpful for recognizing CAPTCHAs, reading images, and identifying menus.

\begin{tcbinfo}{AI Image Description Support by System}
\begin{itemize}
  \item \textbf{iOS 15+}: Enable "Image Descriptions" in [Settings] > [Accessibility] > [VoiceOver] > [Speech Descriptions]. The system calls a\textbf{local AI model}to automatically generate text descriptions of images (requires iPhone 12 or later).
  \item \textbf{Newer TalkBack app}: Google TalkBack integrates\textbf{image Q\& A}in newer versions—you can ask "What's in this image?" and AI answers in real time (not strictly dependent on Android system version—just update the TalkBack app to experience it).
  \item \textbf{Third-party apps}: "Seeing AI" (iOS), "Lookout" (Android), "Xiaoyi Bangbang" (Huawei) are all excellent AI visual assistants that can recognize objects, read documents, and identify facial expressions.
\end{itemize}

\end{tcbinfo}

\begin{tcbpractice}{Use Cases}
Try asking about a street photo, “How many people are in this image? What are they doing?”, or open a food delivery app and let AI read out the dish names in the menu image—\textbf{this is true “tech making life more equal”}.

\end{tcbpractice}

\booksubsection{“Ears” for Hearing-Impaired Users—Complete Live Captions Guide}

Live Captions is a\textbf{system-level speech-to-text}feature—especially useful for hearing-impaired users in noisy environments, libraries, or when communicating in a foreign language.

\begin{tcbinfo}{Enabling and Optimizing Live Captions}
\textbf{iOS}: [Settings] > [Accessibility] > [Live Captions], turn on the switch. The system detects all device audio (calls, videos, podcasts, FaceTime, voice memos) and automatically displays captions at the top of the screen.

\textbf{Android 12+}: Press the\textbf{volume keys}tap the “Live Caption” icon in the popup; or enable in [Settings] > [Accessibility] > [Live Caption].

\textbf{Accuracy optimization tips}:

\begin{enumerate}
  \item Use in a\textbf{quiet environment}as much as possible, away from air conditioners and machinery noise;
  \item The speaker\textbf{articulates clearly and speaks at a moderate pace}for the highest recognition accuracy;
  \item Multi-person conversations reduce accuracy—use "show only one person" mode;
  \item Supports mainstream languages like Mandarin Chinese, Cantonese, and English;\textbf{minority languages have lower recognition rates};
  \item Background music or voice interference during speech will cause intermittent recognition.
\end{enumerate}

\end{tcbinfo}

\booksubsection{Dual-Channel Call Scenarios: Speech-to-Text + Text-to-Speech}

When users\textbf{can neither hear nor speak}(such as sudden voice loss, laryngeal surgery recovery, severe stuttering), the phone can still act as a “spokesperson.”

\begin{tcbinfo}{Real-Time Call Transcription (TTY / RTT)}
\textbf{iOS}: [Settings] > [Accessibility] > [RTT/TTY], after enabling, a keyboard appears on screen during calls—the other party's speech is transcribed to text in real time, and\textbf{your typed text is also read aloud}to them.

\textbf{Android}: [Settings] > [Accessibility] > [Live Transcribe] + [Sound Notifications]. “Live Transcribe” displays the other party's speech as text in real time; “Sound Notifications” converts your input to speech when voice isn't convenient.

\textbf{Compatibility}:

\begin{itemize}
  \item \textbf{RTT (Real-Time Text)}: Both parties need support; text appears as you type with no delay; mainstream carriers now support it;
  \item \textbf{TTY (Teletypewriter)}: Requires a dedicated TTY device; it is gradually being replaced by RTT;
  \item WeChat/QQ\textbf{voice-to-text feature}is also an alternative: long-press a voice message "Convert to Text."
\end{itemize}

\end{tcbinfo}

\booksubsection{“Small but Practical” Features for Everyone}

\begin{itemize}
  \item \textbf{Back Tap}(iOS 14+): [Settings] > [Accessibility] > [Touch] > [Back Tap], can bind\textbf{double-tap/triple-tap}to trigger screenshot, lock screen, control center, etc. Especially handy when carrying packages with both hands.
  \item \textbf{Assistive Touch / Floating Ball}: iOS [Accessibility] > [Touch] > [Assistive Touch], Android [Accessibility] > [Floating Navigation]. Virtual buttons replace gestures or physical buttons—lifesaver when buttons are unresponsive or for one-handed use.
  \item \textbf{Switch Control}: Control the phone through external switches, single clicks, or head movements—\textbf{designed for users with physical disabilities}. iOS [Accessibility] > [Switch Control], Android [Accessibility] > [Switch Access].
  \item \textbf{Sound Recognition}(iOS): [Accessibility] > [Sound Recognition], after enabling, the phone "listens" for doorbells, baby cries, fire alarms, sirens, etc., and\textbf{upon recognition, vibrates + sends a screen notification}—very valuable for hearing-impaired users and during sleep.
  \item \textbf{LED Flash Alerts}: iOS [Accessibility] > [Audio/Visual] > [LED Flash for Alerts], Android [Settings] > [Notifications] > [Flash Alerts]. The flash blinks for incoming calls and notifications—\textbf{sensible even without sound in noisy environments}.
  \item \textbf{Mono Audio}: [Settings] > [Accessibility] > [Audio] > [Mono Audio], merges left and right channels. Users with single-ear hearing loss won't lose half the audio information.
\end{itemize}

\begin{tcbinfo}{Apple's Accessibility Design Philosophy}
Apple invests heavily in accessibility, updating and adding assistive features every year. Apple’s website has a dedicated accessibility page detailing how to use each feature. iOS accessibility is typically more complete and more polished than Android’s, which is a major reason many users with disabilities choose Apple.

\end{tcbinfo}

\begin{tcbpractice}{Spend 30 Seconds to Experience}
Find some free time and do the following experiences:

\begin{enumerate}
  \item Enable\textbf{VoiceOver / TalkBack}and try browsing WeChat with three-finger swipes without using your eyes;
  \item Turn on\textbf{Live Captions}and speak to your phone to check recognition accuracy;
  \item Use\textbf{three-finger tap}to zoom in on text to 200\%—experience the daily life of visually impaired users;
  \item Turn off all the above features and reflect on "if I permanently needed these features, how would I want them designed."
\end{enumerate}

\end{tcbpractice}
